\begin{table}[H]
\centering
\caption{Nama Daerah Perbandingan Klaster Biplot vs Kategori IPM (UNDP) (2024)} \label{tab:app1crossname2024}
\begin{adjustbox}{max width=\textwidth}
\begin{tabular}{|>{\centering\arraybackslash}p{4cm}|p{10cm}|}
\hline
\multicolumn{1}{|c|}{\textbf{Kategori}} & \multicolumn{1}{c|}{\textbf{Daftar Wilayah}} \\
\hline
Cl. 1 (Rendah) \& Tinggi & Ciamis \\
\hline
Cl. 1 (Rendah) \& Sedang & Cianjur, Tasikmalaya, Majalengka, Pangandaran, Pandeglang, Lebak, Serang \\
\hline
Cl. 2 (Sedang) \& Tinggi & Kep. Seribu, Bogor, Bandung, Kuningan, Cirebon, Sumedang, Subang, Purwakarta, Karawang, Bandung Barat, Kota Cirebon, Kota Tasikmalaya, Kota Banjar, Tangerang, Kota Cilegon, Kota Serang \\
\hline
Cl. 2 (Sedang) \& Sedang & Sukabumi, Garut, Indramayu \\
\hline
Cl. 3 (Tinggi) \& Sangat Tinggi & Kota Jakarta Selatan, Kota Jakarta Timur, Kota Jakarta Pusat, Kota Jakarta Barat, Kota Jakarta Utara, Kota Bandung, Kota Bekasi, Kota Depok, Kota Cimahi, Kota Tangerang, Kota Tangerang Selatan \\
\hline
Cl. 3 (Tinggi) \& Tinggi & Bekasi, Kota Bogor, Kota Sukabumi \\
\hline
\end{tabular}
\end{adjustbox}
\end{table}